% !Mode:: "TeX:UTF-8"
% 《玩转 Magel 语言》导言区：文档类选项、文档信息、代码样式、提示盒。
% 由 book.tex 与 tools/check_chapters.py 的试编译壳共同引用。

\usepackage{tabularx}
\pdfstringdefDisableCommands{\def\allowbreak{}}
\setlength{\emergencystretch}{3em}
\usetikzlibrary{arrows.meta,backgrounds,decorations.pathreplacing,fit,positioning,shapes.geometric}

\newif\ifrelease
% 正式发布时取消下面一行的注释，并将另一行注释掉：
% \releasetrue
% 非正式版本（草稿/预览）使用下面一行：
\releasefalse

% ------------------------- 文档信息 -------------------------
\title{玩转 Magel 语言}
\subtitle{一本从零开始、一章只讲一件事的 Magel 入门书}
\author{小兔糖糖OvO}
\institute{真正的高手不会受到环境的影响}
\extrainfo{初次见面，我已经关注你很久了。别看我像个孩子，其实我比任何一位大人都了解这个世界。所以我可以用我的知识，换取你路上的见闻吗？——《原神》纳西妲}

\date{Magel 1.0.0 · 2026年10月15日}
\cover{figure/cover.pdf}

% ------------------------- 代码样式 -------------------------
% magel 源码：带行号、关键字着色
\lstdefinestyle{magel}{
  language=C,
  basicstyle=\ttfamily\small,
  numbers=left,
  numberstyle=\color{black},
  keywordstyle=\color{red!40!orange},
  commentstyle=\color{gray},
  stringstyle=\color[RGB]{128,0,0},
  backgroundcolor=\color{gray!5},
  frame=tb, framerule=1pt,
  breaklines=true, columns=fullflexible, keepspaces=true,
  showstringspaces=false, tabsize=4,
  morecomment=[l]{//},
  morecomment=[s]{/*}{*/},
  morestring=[b]",
  morekeywords={int8,int16,int32,int64,uint8,uint16,uint32,uint64,float32,float64,bool,char,string,String,auto,const,let,struct,enum,interface,generic,fn,match,if,elif,else,while,for,loop,in,return,break,continue,module,import,as,new,delete,move,copy,mut,ref,self,static,true,false,void,abi,annotation,Option,Result,Some,None,Ok,Err,List,Map,Set,Pair,Any,Drop,Closeable,main,println,print,assert,defer,scope,exit,type,yield}
}
% 程序输出
\lstdefinestyle{out}{
  language=,
  basicstyle=\ttfamily\small,
  numbers=none,
  backgroundcolor=\color{gray!5},
  frame=tb, framerule=1pt,
  breaklines=true, columns=fullflexible, keepspaces=true,
  showstringspaces=false,
}
% 命令行
\lstdefinestyle{sh}{
  language=bash,
  basicstyle=\ttfamily\small,
  numbers=none,
  keywordstyle=\color{red!40!orange},
  backgroundcolor=\color{gray!5},
  frame=tb, framerule=1pt,
  breaklines=true, columns=fullflexible, keepspaces=true,
  showstringspaces=false,
}
% 配置文件（toml / json）
\lstdefinestyle{cfg}{
  language=,
  basicstyle=\ttfamily\small,
  numbers=none,
  backgroundcolor=\color{gray!5},
  frame=tb, framerule=1pt,
  breaklines=true, columns=fullflexible, keepspaces=true,
  showstringspaces=false,
}

% ------------------------- 提示盒（模板 anybox 的书稿入口） -------------------------
\newenvironment{tip}[1]{\smallskip\begin{anybox}{#1}}{\end{anybox}}
